%% ========================================================================
%% Appendix C: Glossary
%% ========================================================================

\chapter{Glossary}
\label{app:glossary}

Lampiran ini berisi definisi istilah-istilah teknis yang digunakan dalam buku ini.

%% ------------------------------------------------------------------------
\section*{A}

\begin{description}
    \item[ACL (Access Control List)] Mekanisme kontrol akses yang lebih granular
          dibanding permission Unix tradisional, memungkinkan pengaturan hak akses
          untuk multiple user dan group.

    \item[API (Application Programming Interface)] Sekumpulan definisi dan protokol
          untuk membangun dan mengintegrasikan software aplikasi.
\end{description}

%% ------------------------------------------------------------------------
\section*{B}

\begin{description}
    \item[Bash (Bourne Again Shell)] Shell dan bahasa pemrograman command-line
          yang merupakan default di kebanyakan distribusi Linux.

    \item[Bootloader] Program yang memuat sistem operasi ke memori saat komputer
          dinyalakan. Contoh: GRUB, LILO.
\end{description}

%% ------------------------------------------------------------------------
\section*{C}

\begin{description}
    \item[CLI (Command Line Interface)] Interface berbasis teks untuk berinteraksi
          dengan sistem operasi menggunakan perintah-perintah.

    \item[Cron] Daemon untuk menjadwalkan eksekusi perintah atau script pada
          waktu tertentu.
\end{description}

%% ------------------------------------------------------------------------
\section*{D}

\begin{description}
    \item[Daemon] Program yang berjalan di background dan tidak memiliki
          interface langsung dengan pengguna.

    \item[Distro (Distribution)] Varian Linux yang dikemas dengan kernel,
          utilities, dan software aplikasi tertentu.
\end{description}

%% [Lanjutkan untuk huruf E-Z...]

%% ------------------------------------------------------------------------
\section*{K}

\begin{description}
    \item[Kernel] Inti dari sistem operasi yang mengelola hardware dan
          menyediakan layanan dasar untuk software lainnya.
\end{description}

%% ------------------------------------------------------------------------
\section*{S}

\begin{description}
    \item[Shell] Program yang menyediakan interface untuk pengguna berinteraksi
          dengan sistem operasi.

    \item[System Call] Interface pemrograman untuk layanan yang disediakan
          oleh kernel sistem operasi.

    \item[Systemd] Sistem init modern yang digunakan untuk bootstrap user space
          dan mengelola layanan sistem.
\end{description}

%% [Tambahkan istilah lainnya...]
